\uuid{Vpbq}
\titre{Équation}
\auteur{bodin ; hanani ; mzari}
\organisation{exo7}
\licence{CC BY-NC-SA 4.0}
\datecreate{2020-01-20}
\niveau{L1}
\module{Analyse}
\chapitre{Fonctions circulaires et hyperboliques inverses}
\sousChapitre{Autre}
\difficulte{3}
\theme{Fonctions usuelles, équation}
\exo7qcmid{lille-1.fonctions-usuelles.023}

\contenu{
\begin{question}
\qclassification{126.00}
\qtags{motcle=équation}
Soit $(E)$ l'équation : $ \cos 2x = \sin x $. on note $\cal{S}$ l'ensemble des solutions de $(E)$. Quelles sont les assertions vraies ?

\begin{answers}
    \good{Le domaine de définition de $(E)$  est  $\Rr$.}

    \good{${\cal{S}} = \{\frac{\pi}{6}  +\frac{ 2k\pi}{3}, \, k\in \Zz \}  $}

    \bad{${\cal{S}} =  \{ -\frac{\pi}{2} +2k\pi, \,  k\in \Zz \} $}

    \bad{ ${\cal{S}} =  \{ \frac{\pi}{6} +2k\pi, \,  k\in \Zz \}$}
\end{answers}
\begin{explanations}
$(E) \Leftrightarrow  \cos 2x =  \cos ( \frac{\pi}{2} -x) \Leftrightarrow      \exists k\in \Zz; \,  2x =  \frac{\pi}{2} -x +2k\pi \, \mbox {ou} \, 2x =  -\frac{\pi}{2} +x +2k\pi  \Leftrightarrow      \exists k\in \Zz;\,  x =  \frac{\pi}{6}  +\frac{ 2k\pi}{3} \, \mbox {ou} \, x =  -\frac{\pi}{2} +2k\pi$.
Mais un $x$ de la forme $x =  -\frac{\pi}{2} +2k\pi$ peut s'écrire sous la forme
$x = \frac{\pi}{6}  +\frac{ 2k'\pi}{3}$.
Donc ${\cal{S}} = \{\frac{\pi}{6}  +\frac{ 2k\pi}{3}, \, k\in \Zz \}$.

\end{explanations}
\end{question}
}
